% One Britannica scan as extracted by Nemotron Parse 3.5. Filled by src/nemotron_challenge/stages/latex.py.
% The TITLE, BODY and SOURCE markers (each wrapped in two percent signs) are replaced. Compile with XeTeX (tectonic).
\documentclass[10pt]{article}
\usepackage[a4paper,margin=18mm]{geometry}
\usepackage{fontspec}
\usepackage{multicol}
\usepackage{graphicx}
\usepackage[hidelinks]{hyperref}
\usepackage[normalem]{ulem}
\setlength{\parindent}{0pt}
\setlength{\parskip}{4pt}
\setlength{\columnsep}{8mm}
\begin{document}
\section*{Historical source: CARPENTRY}
\begin{multicols}{2}
\subsection*{\textit{continued:} CARPENTRY}
{\small\itshape [Margin: Currently, | Origin of the Strain on a Tbe beam.]}\par
\sout{As,} \textbf{Ao,} and Ag. But because \sout{A/and-2B} \textbf{Af and eE} are equal and parallel, and \sout{Ex} \textbf{Er} and \sout{f/} \textbf{fi} are also \sout{parallel;} \textbf{parallel,} as also \sout{or} \textbf{er} and \sout{jp,} \textbf{fp,} it is evident, that \sout{if/s} \textbf{if is} equal to \sout{eR,} \textbf{rE,} or to \sout{F,} \textbf{oF,} and \sout{sA} \textbf{iA} is \sout{used} \textbf{equal} to \sout{rG,} \textbf{re,} or to \sout{Gg?.} \textbf{Gg.} Therefore the \sout{forefore Ags Aso,} \textbf{four forces Ag, Ao,} Ap, \sout{Av,} \textbf{Ai,} are equal to \sout{AC} \textbf{AG} and \sout{AC.} \textbf{AF.} Therefore \sout{Ac} \textbf{AG} is the compression of the beam AD, or \textbf{the force pressing it on D, and AF is} the force pressing it on the beam HD. The proportion of these pressures, therefore, is not affected by the form of the \sout{point.} \textbf{joint.} This remark is important: for many carpenters think the form and direction of the \sout{hitting} \textbf{butting} joint of \sout{greater} \textbf{great} importance; and even the \sout{checks} \textbf{theorist,} by not prosecuting the general principle through all its consequences, may be led into an \sout{order.} \textbf{error.} The form of the joint is of no importance, in as far as it affects the strains in the direction of the \sout{beams,} \textbf{beams;} but it is often of \sout{pure} \textbf{great} consequence, in respect to its own firmness, and the effect it may have \sout{to brushing} \textbf{in bruising} the piece on which it acts, or being crippled by it. The same \sout{component} \textbf{compression} of AB, and the same thrust on the point D by the \sout{intersection} \textbf{intervention} of AD, will obtain, in whatever way the original pressure on the end A is produced. Thus, \sout{supporting} \textbf{supposing} that a cord is made fast at \sout{all,} \textbf{A,} and \sout{parallel} \textbf{pulled} in the direction AE, and with the same force, the beam AD will be equally compressed, and the prop D must react with the same \sout{forces.} \textbf{force.} But, it often happens \sout{to} \textbf{that} the obliquity of the pressure on AD, instead of \sout{expressing} \textbf{compressing} it, stretches it; and we desire to know what tension it sustains. \sout{Or they} \textbf{Of this} we have a familiar example in a common roof. Let the two \sout{others} \textbf{rafters} AC, \sout{(D) (the 18),} \textbf{AD (fig. 12.),} press on the tie-beam DC. We may suppose \sout{this} \textbf{the} whole weight to \sout{B was very length between} \textbf{press vertically on} the ridge A, as if \sout{it would} \textbf{a weight} B were \sout{hump} \textbf{hung} on there. (See Note \sout{CC).} \textbf{CC.)} We may represent this weight by the \sout{position A} \textbf{portion Ab} of the vertical or \sout{plane lines, intersect} \textbf{plumb line, intercepted} between the ridge and the beam. Then drawing bf and \sout{agr} \textbf{bg} parallel to AD and AC, Ag and \sout{Auf} \textbf{Af} will \sout{represented} \textbf{represent} the \sout{pressure} \textbf{pressures} on AC and AD. \sout{Preductor} \textbf{Produce} AC till CH be equal to Af. The point C is forced out in this direction, and with a force \sout{represents} \textbf{represented} by \sout{the} \textbf{this} line. As this \sout{from one} \textbf{force} is not perpendicularly across the beam, it evidently stretches it; and this extending force must be \sout{withdrawn byan} \textbf{withstood by an} equal force pulling it in the opposite direction. This must arise \sout{than} \textbf{from} a similar \sout{colour thread} \textbf{oblique thrust} of the opposite \sout{rather} \textbf{rafter} on the other end D. We concern ourselves only with this extension at present; but we see that the cohesion of the beam does nothing but supply the balance to the extending forces. It must still be supported externally, that it \sout{\_any resist\_,} \textbf{may resist,} and by \sout{relying obligibly,} \textbf{resisting obliquely,} be stretched. The points C and D are supported on the \sout{walls. This gives that} \textbf{walls, which they press in} the directions CK and DO, parallel to \sout{AA.} \textbf{Ab.} If we draw HK parallel to DC, and HI parallel to CK (that is to Ab), \sout{meaning Dr.} \textbf{meeting DC} produced in \sout{it,} \textbf{I,} it follows from the composition of forces, that the point C would be supported by the two forces KC and IC. In \sout{his} \textbf{like} manner, making DN = \sout{A, g,} \textbf{Ag,} and completing the parallelogram \sout{DMND,} \textbf{DMNO,} the point D would be supported by the forces OD and MD. If we draw go \sout{and/k} \textbf{and fk} parallel to DC, it is plain that they are equal to NO and CI, while Ao and Ak are equal to DO and CK, and \sout{Ak} \textbf{Ab} is equal to the sum of DO and Carpentry. CK \sout{because} \textbf{(because} it is equal to Ao + \sout{A4).} \textbf{Ak).} The \sout{weights} \textbf{weight} of the roof is equal to its vertical pressure on the \sout{wall.} \textbf{walls.} Thus we see, that while a pressure on A, in the direction AB, produces the strains AF and AG, on the pieces AC and AD, it has excited a strain Cl or DM in the place DC. And this completes the mechanism of a frame; for all derive their efficacy from the triangles of which they are composed, as will \sout{ape your} \textbf{appear} more clearly as we proceed. DO : DM : OK : DI : CM : = OO A Where Do : CK - CE: All \sout{\_Or,} \textbf{Or,} the pressures on the \sout{points\_} \textbf{points} C \sout{\_and} \textbf{and} D, in the \sout{direc tion} \textbf{direction} of the \sout{training} \textbf{straining} force \sout{A, or the increally proportioned} \textbf{Ab, are reciprocally proportional} to the portions of \sout{DE Intercepted By Al.\_} \textbf{DC intercepted by Ab.} Also, since A is = DO + CK, we have A\textbackslash{}; cK - CE « + d (or CD) : 50, and A\textbackslash{}; DO = DO - dz. In general, any two of the three parallel forces AB, DO, CK, are to each other in the reciprocal proportion of the parts of \sout{OD, interconnect and} \textbf{CD, intercepted between} their \sout{direction} \textbf{directions} and the direction of the third. And this explains a still more important office of the frame aDC. I can of the grains, such as D, be supported, an external power acting at A, in the direction AB, and with an intensity which may be measured by AB, may be set in equilibrium, with another acting at C, in the direction CLs, opposite to CK or A, and with an intensity represented by CK, for since the pressure CH is partly attributed by the force IC, or the firmness of the beam DC support as D, the force K, will complete the balance. When you are not attend to the support at D, we conceive the force AB to be balanced by KC, or KC to be balanced by AK. And, in like manner, we may neglect the support or force acting at A, and consider the force DO be balanced by CK. Thus our frame becomes a lever, and we are able \sout{it becomes} to trace the interior mechanical procedure which a \sout{Leres, given} \textbf{Lever gives it its efficacy:} it is \sout{efficacy in a} by the \sout{interaction} \textbf{intervention} of the forces of cohesion, which \sout{cannot} \textbf{connect} the points to which the external forces are applied with the supported point or \sout{flucture,} \textbf{fulcrum,} and with each other. These strains or \sout{pressure AB,} \textbf{pressures Aß,} DO, and CK, not being in the directions of the beams, may be called \sout{\_transverse\_.} \textbf{transverse.} We see that by their means a frame of
\end{multicols}
\section*{Source record}
\begin{itemize}
\item Page ID: \texttt{britannica:309f765b1ecc83e0a7f7762e44b4a927a7555443:train:1680}
\item Scan: \texttt{page-1680-fbb4a5fde0c2.jpg}, 600 x 792 px
\item Image SHA-256: \texttt{\small fbb4a5fde0c2b233b6df827e5998ef2fc15cd2bc73b4eb99b93d91db83199d1b}
\item Extraction: \texttt{nvidia/nemotron-3.5-parse}, raw reply \texttt{artifacts/extract/raw/1680.json}
\item Text is Parse OCR output, normalized (long s as s, line-end hyphens joined). It is not corrected by hand.
\item \sout{Struck-through} words are Parse's reading; \textbf{bold} words are Super VL's correction, checked against the scan image. Both are model output.
\end{itemize}
\end{document}
